%%% Hlavní soubor. Zde se definují základní parametry a odkazuje se na ostatní části. %%%

%% Verze pro jednostranný tisk:
% Okraje: levý 40mm, pravý 25mm, horní a dolní 25mm
% (ale pozor, LaTeX si sám přidává 1in)
\documentclass[12pt,a4paper]{report}
\setlength\textwidth{145mm}
\setlength\textheight{247mm}
\setlength\oddsidemargin{15mm}
\setlength\evensidemargin{15mm}
\setlength\topmargin{0mm}
\setlength\headsep{0mm}
\setlength\headheight{0mm}
% \openright zařídí, aby následující text začínal na pravé straně knihy
\let\openright=\clearpage

%% Pokud tiskneme oboustranně:
% \documentclass[12pt,a4paper,twoside,openright]{report}
% \setlength\textwidth{145mm}
% \setlength\textheight{247mm}
% \setlength\oddsidemargin{14.2mm}
% \setlength\evensidemargin{0mm}
% \setlength\topmargin{0mm}
% \setlength\headsep{0mm}
% \setlength\headheight{0mm}
% \let\openright=\cleardoublepage

%% Vytváříme PDF/A-2u
\usepackage[a-2u]{pdfx}

%% Přepneme na českou sazbu a fonty Latin Modern
\usepackage[czech]{babel}
\usepackage{lmodern}
\usepackage[T1]{fontenc}
\usepackage{textcomp}

%% Použité kódování znaků: obvykle latin2, cp1250 nebo utf8:
\usepackage[utf8]{inputenc}

%%% Další užitečné balíčky (jsou součástí běžných distribucí LaTeXu)
\usepackage{amsmath}        % rozšíření pro sazbu matematiky
\usepackage{amsfonts}       % matematické fonty
\usepackage{amsthm}         % sazba vět, definic apod.
\usepackage{bbding}         % balíček s nejrůznějšími symboly
			    % (čtverečky, hvězdičky, tužtičky, nůžtičky, ...)
\usepackage{bm}             % tučné symboly (příkaz \bm)
\usepackage{graphicx}       % vkládání obrázků
\usepackage{fancyvrb}       % vylepšené prostředí pro strojové písmo
\usepackage{indentfirst}    % zavede odsazení 1. odstavce kapitoly
\usepackage{natbib}         % zajištuje možnost odkazovat na literaturu
			    % stylem AUTOR (ROK), resp. AUTOR [ČÍSLO]
\usepackage[nottoc]{tocbibind} % zajistí přidání seznamu literatury,
                            % obrázků a tabulek do obsahu
\usepackage{icomma}         % inteligetní čárka v matematickém módu
\usepackage{dcolumn}        % lepší zarovnání sloupců v tabulkách
\usepackage{booktabs}       % lepší vodorovné linky v tabulkách
\usepackage{paralist}       % lepší enumerate a itemize
\usepackage{xcolor}         % barevná sazba

\usepackage[colorinlistoftodos]{todonotes}
\usepackage{csvsimple}

%%% Údaje o práci

% Název práce v jazyce práce (přesně podle zadání)
\def\NazevPrace{Umělá inteligence pro hru Calico}

% Název práce v angličtině
\def\NazevPraceEN{Artificial Intelligence for the Calico Game}

% Jméno autora
\def\AutorPrace{Dita Chabičovská}

% Rok odevzdání
\def\RokOdevzdani{2024}

% Název katedry nebo ústavu, kde byla práce oficiálně zadána
% (dle Organizační struktury MFF UK, případně plný název pracoviště mimo MFF)
\def\Katedra{Katedra teoretické informatiky a matematické logiky}
\def\KatedraEN{Department of Theoretical Computer Sci\-ence and Mathematical Logic}

% Jedná se o katedru (department) nebo o ústav (institute)?
\def\TypPracoviste{Katedra}
\def\TypPracovisteEN{Department}

% Vedoucí práce: Jméno a příjmení s~tituly
\def\Vedouci{Mgr. Martin Pilát, Ph.D.}

% Pracoviště vedoucího (opět dle Organizační struktury MFF)
\def\KatedraVedouciho{Katedra teoretické informatiky a matematické logiky}
\def\KatedraVedoucihoEN{Department of Theoretical Computer Sci\-ence and Mathematical Logic}

% Studijní program a obor
\def\StudijniProgram{Informatika}
\def\StudijniObor{Obecná informatika}

% Nepovinné poděkování (vedoucímu práce, konzultantovi, tomu, kdo
% zapůjčil software, literaturu apod.)
\def\Podekovani{%
Poděkování.
}

% Abstrakt (doporučený rozsah cca 80-200 slov; nejedná se o zadání práce)
\def\Abstrakt{%
Abstrakt.
}
\def\AbstraktEN{%
Abstract.
}

% 3 až 5 klíčových slov (doporučeno), každé uzavřeno ve složených závorkách
\def\KlicovaSlova{%
{Calico} {hra} {umělá inteligence}
}
\def\KlicovaSlovaEN{%
{Calico} {board game} {artificial intelligence}
}

%% Balíček hyperref, kterým jdou vyrábět klikací odkazy v PDF,
%% ale hlavně ho používáme k uložení metadat do PDF (včetně obsahu).
%% Většinu nastavítek přednastaví balíček pdfx.
\hypersetup{unicode}
\hypersetup{breaklinks=true}

%% Definice různých užitečných maker (viz popis uvnitř souboru)
%%% Tento soubor obsahuje definice různých užitečných maker a prostředí %%%
%%% Další makra připisujte sem, ať nepřekáží v ostatních souborech.     %%%

%%% Drobné úpravy stylu

% Tato makra přesvědčují mírně ošklivým trikem LaTeX, aby hlavičky kapitol
% sázel příčetněji a nevynechával nad nimi spoustu místa. Směle ignorujte.
\makeatletter
\def\@makechapterhead#1{
  {\parindent \z@ \raggedright \normalfont
   \Huge\bfseries \thechapter. #1
   \par\nobreak
   \vskip 20\p@
}}
\def\@makeschapterhead#1{
  {\parindent \z@ \raggedright \normalfont
   \Huge\bfseries #1
   \par\nobreak
   \vskip 20\p@
}}
\makeatother

% Toto makro definuje kapitolu, která není očíslovaná, ale je uvedena v obsahu.
\def\chapwithtoc#1{
\chapter*{#1}
\addcontentsline{toc}{chapter}{#1}
}

% Trochu volnější nastavení dělení slov, než je default.
\lefthyphenmin=2
\righthyphenmin=2

% Zapne černé "slimáky" na koncích řádků, které přetekly, abychom si
% jich lépe všimli.
\overfullrule=1mm

%%% Makra pro definice, věty, tvrzení, příklady, ... (vyžaduje baliček amsthm)

\theoremstyle{plain}
\newtheorem{veta}{Věta}
\newtheorem{lemma}[veta]{Lemma}
\newtheorem{tvrz}[veta]{Tvrzení}

\theoremstyle{plain}
\newtheorem{definice}{Definice}

\theoremstyle{remark}
\newtheorem*{dusl}{Důsledek}
\newtheorem*{pozn}{Poznámka}
\newtheorem*{prikl}{Příklad}

%%% Prostředí pro důkazy

\newenvironment{dukaz}{
  \par\medskip\noindent
  \textit{Důkaz}.
}{
\newline
\rightline{$\qedsymbol$}
}

%%% Prostředí pro sazbu kódu, případně vstupu/výstupu počítačových
%%% programů. (Vyžaduje balíček fancyvrb -- fancy verbatim.)

\DefineVerbatimEnvironment{code}{Verbatim}{fontsize=\small, frame=single}

%%% Prostor reálných, resp. přirozených čísel
\newcommand{\R}{\mathbb{R}}
\newcommand{\N}{\mathbb{N}}

%%% Užitečné operátory pro statistiku a pravděpodobnost
\DeclareMathOperator{\pr}{\textsf{P}}
\DeclareMathOperator{\E}{\textsf{E}\,}
\DeclareMathOperator{\var}{\textrm{var}}
\DeclareMathOperator{\sd}{\textrm{sd}}

%%% Příkaz pro transpozici vektoru/matice
\newcommand{\T}[1]{#1^\top}

%%% Vychytávky pro matematiku
\newcommand{\goto}{\rightarrow}
\newcommand{\gotop}{\stackrel{P}{\longrightarrow}}
\newcommand{\maon}[1]{o(n^{#1})}
\newcommand{\abs}[1]{\left|{#1}\right|}
\newcommand{\dint}{\int_0^\tau\!\!\int_0^\tau}
\newcommand{\isqr}[1]{\frac{1}{\sqrt{#1}}}

%%% Vychytávky pro tabulky
\newcommand{\pulrad}[1]{\raisebox{1.5ex}[0pt]{#1}}
\newcommand{\mc}[1]{\multicolumn{1}{c}{#1}}


%% Titulní strana a různé povinné informační strany
\begin{document}
%%% Titulní strana práce a další povinné informační strany

%%% Titulní strana práce

\pagestyle{empty}
\hypersetup{pageanchor=false}

\begin{center}

\centerline{\mbox{\includegraphics[width=166mm]{../img/logo-cs.pdf}}}

\vspace{-8mm}
\vfill

{\bf\Large BAKALÁŘSKÁ PRÁCE}

\vfill

{\LARGE\AutorPrace}

\vspace{15mm}

{\LARGE\bfseries\NazevPrace}

\vfill

\Katedra

\vfill

{
\centerline{\vbox{\halign{\hbox to 0.45\hsize{\hfil #}&\hskip 0.5em\parbox[t]{0.45\hsize}{\raggedright #}\cr
Vedoucí bakalářské práce:&\Vedouci \cr
\noalign{\vspace{2mm}}
Studijní program:&\StudijniProgram \cr
\noalign{\vspace{2mm}}
Studijní obor:&\StudijniObor \cr
}}}}

\vfill

% Zde doplňte rok
Praha \RokOdevzdani

\end{center}

\newpage

%%% Následuje vevázaný list -- kopie podepsaného "Zadání bakalářské práce".
%%% Toto zadání NENÍ součástí elektronické verze práce, nescanovat.

%%% Strana s čestným prohlášením k bakalářské práci

\openright
\hypersetup{pageanchor=true}
\pagestyle{plain}
\pagenumbering{roman}
\vglue 0pt plus 1fill

\noindent
Prohlašuji, že jsem tuto bakalářskou práci vypracoval(a) samostatně a výhradně
s~použitím citovaných pramenů, literatury a dalších odborných zdrojů.
Tato práce nebyla využita k získání jiného nebo stejného titulu.

\medskip\noindent
Beru na~vědomí, že se na moji práci vztahují práva a povinnosti vyplývající
ze zákona č. 121/2000 Sb., autorského zákona v~platném znění, zejména skutečnost,
že Univerzita Karlova má právo na~uzavření licenční smlouvy o~užití této
práce jako školního díla podle §60 odst. 1 autorského zákona.

\vspace{10mm}

\hbox{\hbox to 0.5\hsize{%
V \hbox to 6em{\dotfill} dne \hbox to 6em{\dotfill}
\hss}\hbox to 0.5\hsize{\dotfill\quad}}
\smallskip
\hbox{\hbox to 0.5\hsize{}\hbox to 0.5\hsize{\hfil Podpis autora\hfil}}

\vspace{20mm}
\newpage

%%% Poděkování

\openright

\noindent
\Podekovani

\newpage

%%% Povinná informační strana bakalářské práce

\openright

\vbox to 0.5\vsize{
\setlength\parindent{0mm}
\setlength\parskip{5mm}

Název práce:
\NazevPrace

Autor:
\AutorPrace

\TypPracoviste:
\Katedra

Vedoucí bakalářské práce:
\Vedouci, \KatedraVedouciho

Abstrakt:
\Abstrakt

Klíčová slova:
\KlicovaSlova

\vss}\nobreak\vbox to 0.49\vsize{
\setlength\parindent{0mm}
\setlength\parskip{5mm}

Title:
\NazevPraceEN

Author:
\AutorPrace

\TypPracovisteEN:
\KatedraEN

Supervisor:
\Vedouci, \KatedraVedoucihoEN

Abstract:
\AbstraktEN

Keywords:
\KlicovaSlovaEN

\vss}

\newpage

\openright
\pagestyle{plain}
\pagenumbering{arabic}
\setcounter{page}{1}


%%% Strana s automaticky generovaným obsahem bakalářské práce

\tableofcontents

\listoftodos

%%% Jednotlivé kapitoly práce jsou pro přehlednost uloženy v samostatných souborech
\chapter*{Úvod}
\addcontentsline{toc}{chapter}{Úvod}

Ok, let's go.
\todo[inline]{
Úvod

existuje webová aplikace a nová hra na Steamu

neexistuje hra proti počítačovámu hráči? (musím zkouknout ten Steam)

jsem první, kdo dělá nějakou analýzu

chci udělat počítačového hráče a zanalyzovat optimální nastavení na začátku
}




%%% Fiktivní kapitola s ukázkami sazby

\chapter{Popis hry}

\section{Základní pravidla}

Hra Calico je založena na principu postupného přikládání dílků. 

Každý hráč má vlastní herní desku s políčky v mřížce velikosti 5x5. Na desku postupně doplňuje dílky s různými barvami a vzory a získává body podle zadaných pravidel na seskupování barev a vzorů v mřížce. Hráči se při vybírání dílků střídají. Hra končí, když mají všichni hráči vyplněné celé herní pole.

Vítězí hráč, kterému se podaří získat nejvíce bodů.

\subsubsection*{Herní materiál???}
6 barev, 6 vzorů

Hra obsahuje 3 dílky od každé kombinace barvy a vzoru (36 kombinací), tedy 108 herních dílků.
\subsubsection*{Začátek hry}
Každý hráč dostane svou hrací desku.

Na stůl se vyloží 3 dílky z pytlíku. 

Náhodně se vybere, ze kterých vzorů budeme skládat trojice/čtveřice/pětice.
\subsubsection*{Průběh hry}

V každém kole hry se všichni hráči vystřídají na tahu ve stejném na začátku stanoveném pořadí. Hráč na tahu vybere 1 ze 3 možných dílků a přidá ho na svou hrací desku na libovolné neobsazené políčko. Poté doplní možnosti použitelných dílků.

\subsubsection*{Podmínky získání bodů}
Za každou trojici sousedících stejnobarevných dílků hráč získává 3 body.

Vzory jsou náhodně rozděleny do 3 dvojic s následujícím bodováním:
\begin{enumerate}
    \item Sousedící trojice se stejným vzorem = 3 body
    \item Sousedící čtveřice se stejným vzorem = 5 bodů
    \item Sousedící pětice se stejným vzorem = 7 bodů
\end{enumerate}

\subsubsection*{Konec hry}

Hra končí když každý hráč naplní celé své hrací pole.


\section{Možná rozšíření}

Tohle je nějaká family verze -> víc casual. Běžně se hraje s úkolovými dílky.

př. těch úkolových dílků

možnost specifikovat tvar skupin dílků se stejným vzorem

víc her za sebou, každá s jiným pravidlem (např. bez zisku knoflíků)

\section{Mé úpravy oproti původním pravidlům}

bez žetonků kočiček a knoflíků, takže nekontroluji víc kočiček vedle sebe, nebo jestli je nový barevný cluster skutečně kompaktní (např 6 dílků = 2 clustery bez ohledu na to, jestli ta druhá trojice sousedí), protože by to bylo náročné na rozhodování v případě sjednocení více clusterů, který z nich chci vzít

automatické počítání bodů v průběhu hry 

body za posbírání knoflíků všech barev ?? (můžu přidat)


\chapter{Analýza hry}

\section{Vlastnosti hry?}
vidím všechno - herní desky ostatních, stejné možnosti pro každého

hra s úplnou informací (alespoň v mojí verzi)

hra s nenulovým součtem

sekvenční hra (hráči se střídají na tahu)

intenzivní množství pravděpodobnosti, což je bída

mám přesně daný počet kol = 22, konečná hra


\section{AI - přístupy ?}
jaké přístupy pro ai dávají a nedávají smysl

basic tree search, monte carlo tree search

evoluce
\subsection{Strom}
\subsection{MonteCarlo treesearch}

\subsection{Evoluční algoritmus}
nějaká teorie = co je evoluční algoritmus a tak
\subsubsection*{Incializace}
\subsubsection*{Selekce}
\subsubsection*{Křížení}
\subsubsection*{Mutace}


\chapter{Implementace agentů}

rozdělení na single a multiplayer

možná evoluce? pak rozdělení na ručně psané a optimalizující idk

monte carlo tree search

rozdělení hry na fáze

\section{Hra pro jednoho hráče}
\section{Hra pro více hráčů}

\chapter{Implementace hry}


\section{Požadavky na program}

\section{UserDoc}

\subsection{Ovládání konzolové aplikace}

Pro spuštění a ovládání hry postupujte podle instrukcí v konzoli.
Všechny instrukce zadávejte ve formě celých čísel, (popř y/n, pokud vás o to instrukce žádá).

Barvy jsou kódovány čísly 1-6, vzory velkými písmeny A-F.
Prázdná políčka jsou označena -{}-.
Políčka označená XX jsou blokovaná a nelze je obsadit.

\subsection{Herní módy}
\subsection{Testování}
\section{DevDoc}
\subsection{Objektový návrh}
\chapter{Výsledky}

\section{Analýza nastavení hry}
\subsection{Výběr úkolových dílků}
zatím testované na agentovi s tou nějakou utility fcí (11)

kombinace 1,2,3 nejhorší -> většinou průměr pod 39,5, max k 39,7

kombinace 4,5,6 nejlepší (obecně dobrá kombinace 4,6) -> kolem 45, občas některá kombinace až 46

nějaká super tabulka na všechny ty kombinace (bez závislosti na pořadí, tzn nějaký průměr)

tabulka na rozdíly kombinací 4,5,6

pořadí výběru tasků (tedy jejich rozložení na desce) dělá rozdíly, ale ne příliš velké
\begin{table}[h!]
\centering
\begin{tabular}{|c|c|}%
    \hline
    \bfseries Kombinace & \bfseries Průměrné skóre% specify table head
    \\\hline
    \csvreader[head to column names, separator=semicolon,
    late after line=\\,
    late after last line=\\\hline]
    {./tabulky/testTasks.csv}{}% use head of csv as column names
    { \tasks & \averageScore}% specify your coloumns here
\end{tabular}
\caption{Porovnání výsledků pro různé kombinace úkolových dílků}
\end{table}

\begin{table}[h!]
\centering
\begin{minipage}{.45\linewidth}
\centering
\begin{tabular}{|c|c|}%
    \hline
    \bfseries Kombinace & \bfseries Průměrné skóre% specify table head
    \\\hline
    \csvreader[head to column names, separator=semicolon,
    late after line=\\,
    late after last line=\\\hline]
    {./tabulky/testTasks123.csv}{}% use head of csv as column names
    {\tasks & \score}% specify your coloumns here
\end{tabular}
\end{minipage}
\begin{minipage}{.45\linewidth}
\centering
\begin{tabular}{|c|c|}%
    \hline
    \bfseries Kombinace & \bfseries Průměrné skóre% specify table head
    \\\hline
    \csvreader[head to column names, separator=semicolon,
    late after line=\\,
    late after last line=\\\hline]
    {./tabulky/testTasks235.csv}{}% use head of csv as column names
    {\tasks & \score}% specify your coloumns here
\end{tabular}
\end{minipage}
\caption{Porovnání různých umístění kombinací s nejhoršími výsledky}
\end{table}

\begin{table}[h!]
\centering
\begin{minipage}{.45\textwidth}
\centering
\begin{tabular}{|c|c|}%
    \hline
    \bfseries Kombinace & \bfseries Průměrné skóre% specify table head
    \\\hline
    \csvreader[head to column names, separator=semicolon,
    late after line=\\,
    late after last line=\\\hline]
    {./tabulky/testTasks146.csv}{}% use head of csv as column names
    {\tasks & \score}% specify your columns here
\end{tabular}
\end{minipage}
\begin{minipage}{.45\textwidth}
\centering
\begin{tabular}{|c|c|}%
    \hline
    \bfseries Kombinace & \bfseries Průměrné skóre% specify table head
    \\\hline
    \csvreader[head to column names, separator=semicolon,
    late after line=\\,
    late after last line=\\\hline]
    {./tabulky/testTasks456.csv}{}% use head of csv as column names
    {\tasks & \score}% specify your columns here
\end{tabular}
\end{minipage}
\caption{Porovnání různých umístění kombinací s nejlepšími výsledky}
\end{table}

\vspace{5pt}

\todo[inline, color=blue!40]{výběr desky? 4 možnosti}

\section{Porovnání výsledků pro hru jednoho hráče}
porovnání v single oproti random a lidskému hráči

\begin{table}[h!]
\centering

\begin{tabular}{|c|c|c|c|c|c|}%
    \hline
    \bfseries Agent & \bfseries Skóre & \bfseries Knoflíky & \bfseries Kočky & \bfseries Max & \bfseries Min
    \\\hline
    \csvreader[head to column names, separator=semicolon,
    late after line=\\,
    late after last line=\\\hline]
    {./tabulky/testAll.csv}{}% use head of csv as column names
    { \agent & \averageScore & \buttons & \cats & \best & \lowest}% specify your coloumns here
\end{tabular}
\caption{Porovnání agentů ve hře pro jednoho hráče}
\end{table}

\section{Porovnání výsledků pro hru více hráčů}
porovnání multi oproti random hráči

\section{Porovnání strategií pro hru jednoho a více hráčů}

porovnání parametrů evoluce??



\chapter*{Závěr}
\addcontentsline{toc}{chapter}{Závěr}

Už to dál nezvládám, skutečně jsem na konci svých možností.
\todo[inline]{Závěr}

%%% Seznam použité literatury
%%% Seznam použité literatury (bibliografie)
%%%
%%% Pro vytváření bibliografie používáme bibTeX. Ten zpracovává
%%% citace v textu (např. makro \cite{...}) a vyhledává k nim literaturu
%%% v souboru literatura.bib.
%%%
%%% Příkaz \bibliographystyle určuje, jakým stylem budou citovány odkazy
%%% v textu. V závorce je název zvoleného souboru .bst. Styly plainnat
%%% a unsrt jsou standardní součástí latexových distribucí. Styl czplainnat
%%% je dodáván s touto šablonou a bibTeX ho hledá v aktuálním adresáři.

\bibliographystyle{czplainnat}    %% Autor (rok) s českými spojkami
% \bibliographystyle{plainnat}    %% Autor (rok) s anglickými spojkami
% \bibliographystyle{unsrt}       %% [číslo]

\renewcommand{\bibname}{Seznam použité literatury}

%%% Vytvoření seznamu literatury. Pozor, pokud jste necitovali ani jednu
%%% položku, seznam se automaticky vynechá.

\bibliography{cs/literatura/literatura}

%%% Kdybyste chtěli bibliografii vytvářet ručně (bez bibTeXu), lze to udělat
%%% následovně. V takovém případě se řiďte normou ISO 690 a zvyklostmi v oboru.

% \begin{thebibliography}{99}
%
% \bibitem{lamport94}
%   {\sc Lamport,} Leslie.
%   \emph{\LaTeX: A Document Preparation System}.
%   2. vydání.
%   Massachusetts: Addison Wesley, 1994.
%   ISBN 0-201-52983-1.
%
% \end{thebibliography}


%%% Obrázky v bakalářské práci
%%% (pokud jich je malé množství, obvykle není třeba seznam uvádět)
\listoffigures

%%% Tabulky v bakalářské práci (opět nemusí být nutné uvádět)
%%% U matematických prací může být lepší přemístit seznam tabulek na začátek práce.
\listoftables

%%% Použité zkratky v bakalářské práci (opět nemusí být nutné uvádět)
%%% U matematických prací může být lepší přemístit seznam zkratek na začátek práce.
\chapwithtoc{Seznam použitých zkratek}

%%% Přílohy k bakalářské práci, existují-li. Každá příloha musí být alespoň jednou
%%% odkazována z vlastního textu práce. Přílohy se číslují.
%%%
%%% Do tištěné verze se spíše hodí přílohy, které lze číst a prohlížet (dodatečné
%%% tabulky a grafy, různé textové doplňky, ukázky výstupů z počítačových programů,
%%% apod.). Do elektronické verze se hodí přílohy, které budou spíše používány
%%% v elektronické podobě než čteny (zdrojové kódy programů, datové soubory,
%%% interaktivní grafy apod.). Elektronické přílohy se nahrávají do SISu a lze
%%% je také do práce vložit na CD/DVD. Povolené formáty souborů specifikuje
%%% opatření rektora č. 72/2017.
\appendix
\chapter{Přílohy}

\section{První příloha}

\openright
\end{document}
